\subsection{Teststrategie}
\label{subsec:teststrategie}

Die Verifikation ordnet jede Beobachtung einem Artefakt, einer ausgeführten
Prüfung und einem commitgebundenen Ergebnis zu. Testcode allein belegt keine
Ausführung, eine Workflow-Datei keinen erfolgreichen CI-Lauf.
Ausführungsergebnisse und Abnahmeprotokoll wiegen daher stärker als Code,
Konfiguration und Dokumentation
\parencite{runeson2009case-study,iso25010-2023}.

Vier Ebenen decken die Baseline ab: Tooling-Tests prüfen CLI, Profile,
Generator, Konfiguration, Release und Repository; Komponententests Schema,
Backend, Frontend und Tauri/Rust. Integrationstests verbinden PostgreSQL~16,
Konfiguration und Alembic. Buildprüfungen erzeugen Web-, Container- oder native
Artefakte \parencite{kleiveist2026ci}.

Der Umfang folgt dem aktiven Profil. \emph{PASS} bezeichnet eine erfolgreiche,
\emph{FAIL} eine ausgeführte und gescheiterte Prüfung. \emph{Nicht verifiziert}
bedeutet fehlende Evidenz, \emph{außerhalb des Scopes} eine bewusste Grenze;
deaktivierte Gegenstände sind nicht anwendbar. Generierung und Strukturprüfung
adressieren A-01 und A-03, die öffentliche Tooling-Schnittstelle A-04 und
gültige sowie abgewiesene Capability-Kombinationen A-07. Abbildung~\ref{fig:verification-model}
ordnet diese Nachweise ein.

\begin{figure}[H]
  \centering
  \resizebox{0.98\textwidth}{!}{%
  \begin{tikzpicture}
    \node[casePrimary,minimum width=45mm] (local) at (-4.2,0)
      {lokal: Entwicklung\\oder Coding Agent};
    \node[casePrimary,minimum width=45mm] (ci) at (4.2,0)
      {GitHub Actions};

    \node[casePrimary,minimum width=52mm] (control) at (0,-1.7)
      {\texttt{tools/control.py}\\öffentlicher Dispatcher};
    \node[casePrimary,minimum width=58mm] (orchestration) at (0,-3.3)
      {\texttt{tools/quality/control.py}\\Quality-Orchestrierung};

    \node[caseBox,minimum width=45mm,minimum height=19mm] (policy) at (-5,-5.3)
      {\textbf{Policy / Repository}\\Konfiguration, Scan,\\Excludes, Exceptions};
    \node[caseBox,minimum width=45mm,minimum height=19mm] (tools) at (0,-5.3)
      {\textbf{Sprachwerkzeuge}\\Ruff, ESLint, TypeScript,\\Prettier, Clippy, Cargo};
    \node[caseBox,minimum width=45mm,minimum height=19mm] (architecture) at (5,-5.3)
      {\textbf{Architekturregeln}\\Backend- und\\Frontend-Importe};

    \node[casePrimary,minimum width=52mm] (report) at (0,-7.3)
      {Text- oder JSON-Bericht};
    \node[caseOptional,minimum width=52mm] (exit) at (0,-8.7)
      {Severity-Auswertung und Exitcode};

    \draw[caseFlow] (local.south) -- (control.north);
    \draw[caseFlow] (ci.south) -- (control.north);
    \draw[caseFlow] (control) -- (orchestration);
    \coordinate (branch) at (0,-4.25);
    \draw[caseFlow] (orchestration.south) -- (branch) -| (policy.north);
    \draw[caseFlow] (branch) -- (tools.north);
    \draw[caseFlow] (branch) -| (architecture.north);
    \draw[caseFlow] (policy.south) |- (report.west);
    \draw[caseFlow] (tools.south) -- (report.north);
    \draw[caseFlow] (architecture.south) |- (report.east);
    \draw[caseFlow] (report) -- (exit);
  \end{tikzpicture}}
  \caption{Quality-Governance als Entwicklungs- und CI-Schicht außerhalb der Produktlaufzeit}
  \label{fig:verification-model}
  \smallskip
  {\footnotesize Quelle: Eigene Darstellung auf Grundlage von
  \textcite{kleiveist2026quality-implementation} und
  \textcite{kleiveist2026quality-ci}.\par}
\end{figure}

\beginnerinput{workspace/statement/600_qualitaet_verifikation/610}
